\question[4]
\begin{parts}
\part Schreibe eine rekursive Funktion \texttt{anzahl\_gerade(a)}, die zurückgibt, wie viele gerade Zahlen in der
Liste \texttt{a} stehen. \texttt{anzahl\_gerade([4, 7, 2, 9, 6])} liefert \texttt{3}.
\part Schreibe eine rekursive Funktion \texttt{maximum(a)}, die die größte Zahl einer nicht leeren Liste
zurückgibt, ohne \texttt{max} zu verwenden. \texttt{maximum([4, 7, 2, 9, 6])} liefert \texttt{9}.
\end{parts}
\begin{solutionbox}{8cm}
\begin{lstlisting}
# a)
def anzahl_gerade(a):
    if len(a) == 0: return 0
    if a[0] % 2 == 0:
        return 1 + anzahl_gerade(a[1:])
    return anzahl_gerade(a[1:])

# b)
def maximum(a):
    if len(a) == 1: return a[0]
    m = maximum(a[1:])          # Maximum des Rests
    if a[0] > m: return a[0]
    return m
\end{lstlisting}
\end{solutionbox}
